% !TeX root = ../main.tex
% 本节数据来自冻结的主方案；运输架次下限试验仅在对照部分使用。
\begingroup
\providecommand{\qThreeFigDir}{figures/q3}
\captionsetup{font=small,skip=4pt}
\setlength{\textfloatsep}{8pt plus 2pt minus 2pt}
\setlength{\floatsep}{6pt plus 2pt minus 2pt}
\setlength{\intextsep}{8pt plus 2pt minus 2pt}
\setlength{\emergencystretch}{1em}
\section{问题三运输与通信中继联合调度}
\label{sec:problem3}

\subsection{问题分析与技术路线}
\subsubsection{运输执行与通信保障的时空耦合}
问题三在前两问的运输能力、逐箱交付和机电资源约束基础上，增加运输无人机在途通信要求。
运输机沿山地航线移动时，其与网关之间的距离和地形遮挡状态不断变化；
即使起点、终点均可通信，也不能据此保证整个飞行和交接过程持续连通。
当直接链路不可用时，需要已到位并完成建链的中继机同时维持接入链路和回传链路。
因此，中继悬停位置、服务时窗和返航安排必须与运输任务一起确定。

本问继承前两问的运输物理模型与加权目标，重新选择货箱组合、机型、访问顺序及执行时刻，
不固定沿用问题二推荐方案的22条运输任务。
以运输候选为基础，为每条物理航迹构造可连续认证的通信方案，
再将所需中继服务区间嵌入CP-SAT模型，联合优化两类任务及资源安排。
其中，运输架次数记为$N_T$，中继架次数记为$N_R$，计入综合目标的总架次数为$N=N_T+N_R$。

\subsubsection{数据整理与通信参数}
运输数据仍包含15个服务区、80个不可拆货箱、8架运输无人机和14组共享电池。
节点、DEM、240条有向航段、逐段剩余载荷、31个有效硬截止以及分段充电规则均与问题二一致。
新增数据包括网关高度、三类链路的收发参数、中继机性能和能源组件库存。
以O01为局部平面坐标原点，位置统一使用米，时刻统一为任务开始后的秒数。
网关G01位于O01上方；中继位置由水平坐标和相对当地DEM的离地高度共同确定。

通信工作频率为2400\,MHz，系统损耗为3\,dB，存在地形遮挡时附加10\,dB损耗，
接收灵敏度与题给安全裕量合成的接收门限为$-90$\,dBm。
可用中继机为2架，能源组件为6组；每组可用能量为3.2\,kWh，返航SOC不得低于20\%。
中继准备、建链、返航后周转时间分别为180\,s、30\,s和300\,s，
悬停与通信功率合计为1.10\,kW，候选离地高度不超过300\,m。

\subsubsection{求解流程}
图~\ref{fig:q3-workflow}给出本问的技术路线。
首先按前两问规则构造运输候选并计算航迹；其次利用地形和双向链路预算，
认证直连、单站中继或分段双站中继方案，提取各站所需服务区间；
随后联合求解运输选择、站内交箱顺序、中继启用与执行时刻；
最后独立复算逐箱交付、机电占用、返航SOC和连续通信。
候选增补与联合排程交替进行，保留通过核验的完整方案。
\begin{figure}[H]
 \centering
 \includegraphics[width=\linewidth]{\qThreeFigDir/q3_01_workflow.pdf}
 \caption{运输候选、连续通信认证与中继联合调度技术路线}
 \label{fig:q3-workflow}
\end{figure}

\subsection{运输与通信中继联合调度模型}
\subsubsection{运输约束的继承与状态表示}
沿用问题二的候选定义$p=(g_p,\mathcal B_p,\pi_p)$，其中$g_p$、$\mathcal B_p$、$\pi_p$
分别表示机型、非空货箱集合和服务区访问顺序，$x_p$表示是否选用该候选。
每条候选重新计算装载、逐段剩余载荷、运输能耗$E_p$、持续时间$D_p$和充电时间，
并满足质量、体积和20\%返航余量要求。所有货箱恰好交付一次：
\begin{equation}
 \sum_{p:\,b\in\mathcal B_p}x_p=1\quad(\forall b),\qquad x_p\in\{0,1\}.
 \label{eq:q3-cover}
\end{equation}
各运输架次仍从O01出发并返回O01，按逐箱交接顺序确定$C_b$，
医疗箱和首批保障箱满足有效硬截止$C_b\le H_b$。
同型无人机占用区间为$[s_p,s_p+D_p)$，
电池占用延续至返航后充满，两类累计占用分别受$(4,2,2)$和$(6,4,4)$的库存限制。

通信计算采用与运输模型相同的三维轨迹。
将爬升、巡航、下降及服务区交接分别表示为分段仿射轨迹$\boldsymbol x_p(t)$，
交接阶段位置固定。每完整航段向上取整至1\,ms造成的尾差作为到站位置的数值等待，
一并检查通信。通信要求覆盖离地至最终返航的全部在途区间$\mathcal T_p$，
地面准备与装载时间仍计入资源占用，但不视为在途通信区间。

\subsubsection{地形遮挡与双向链路预算}
设$u,v$为两个通信端点，三维距离为$d_{uv}$，单位为km。
采用题给自由空间损耗和遮挡附加损耗模型：
\begin{equation}
 L_{uv}=32.45+20\log_{10}f+20\log_{10}d_{uv}
             +10I_{uv}^{\mathrm{blk}},\qquad f=2400\ \mathrm{MHz}.
 \label{eq:q3-pathloss}
\end{equation}
其中$I_{uv}^{\mathrm{blk}}$表示地形遮挡。
对无线电视线经过的闭DEM像元，比较像元高程与视线在该像元内的最低海拔；
地形高于或接触视线时判为遮挡，像元边界及角点接触亦纳入检查。
遮挡通过附加损耗影响链路预算，是否可用由最终接收功率确定。

以$P_u,G_u$分别表示端点$u$的发射功率和天线增益，
$P_{\mathrm{th}}=-90$\,dBm为含安全裕量的门限。
考虑上下行中较弱的一侧，定义双向链路余量
\begin{equation}
 M_{uv}=\min(P_u,P_v)+G_u+G_v-3-L_{uv}-P_{\mathrm{th}},\qquad
 A_{uv}=\mathbf 1_{\{M_{uv}\ge0\}}.
 \label{eq:q3-linkmargin}
\end{equation}
整理题给收发参数后，运输机--网关、运输机--中继接入端以及中继回传端--网关
所允许的最大路径损耗分别为122、116和126\,dB。
三类链路分别计算，不将接入链路可用等同于中继已能连接网关。

\subsubsection{全在途通信约束与实际服务窗口}
设$u_r\in\{0,1\}$表示是否启用候选中继任务$r$，
其悬停位置为$\boldsymbol q_r$，完成建链后的服务窗口为$[a_r,b_r]$。
中继有效状态为$B_r(t)=u_r\mathbf1_{\{a_r\le t\le b_r\}}$。
每个选中运输任务在任意在途时刻均须满足
\begin{equation}
 A_{pG}(t)=1\quad\text{或}\quad
 \exists r:\ B_r(t)A_{pr}(t)A_{rG}=1,
 \qquad \forall t\in\mathcal T_p.
 \label{eq:q3-continuous}
\end{equation}
当直连可用时优先使用G01；需要中继时，接入和回传链路必须同时可用。
每个时刻最多经过一个中继节点，禁止中继之间多跳转发。
同一运输架次可在不同区段使用不同的实际中继，但切换后的中继必须已经到位并处于服务窗口内。

为将连续条件嵌入离散调度，对运输候选$p$构造有限通信选项集合$\mathcal O_p$。
选项$o$规定各连续区段使用G01或某个中继，并为所需站$r$生成相对准备开始时刻的
首末需求偏移$\ell_{por},h_{por}$。引入$y_{po}\in\{0,1\}$，满足
\begin{align}
 \sum_{o\in\mathcal O_p}y_{po}&=x_p, \label{eq:q3-option}\\
 y_{po}=1&\ \Longrightarrow\ 
 u_r=1,\quad a_r\le s_p+\ell_{por},\quad
 s_p+h_{por}\le b_r\quad(\forall r\text{被选项}o\text{使用}).
 \label{eq:q3-servicewindow}
\end{align}
首末偏移形成服务需求包络，中继在包络内连续服务；
毫秒建模时对需求起点向下取整、终点向上取整，防止数值取整缩短保障区间。

\subsubsection{中继时序与能量资源约束}
记$c_r$为中继任务准备开始时刻，$t_r^{out}$、$t_r^{back}$为往返飞行时间，
$d_r=b_r-a_r$为有效服务时长。启用中继任务后，时序满足
\begin{equation}
 a_r=c_r+180+t_r^{out}+30,\qquad
 b_r=a_r+d_r,\qquad R_r=b_r+t_r^{back}.
 \label{eq:q3-relaytime}
\end{equation}
时长在实际求解中保守量化到毫秒。
中继飞行仍采用地形约束下的爬升、水平巡航与下降过程，
巡航海拔至少达到沿程最高地形加50\,m，并不低于起终点海拔。
以$d_r^{xy}$为单程水平距离，$h_r^+,h_r^-$分别为去程爬升和下降高度，
往返运输能耗及任务总能耗写为
\begin{align}
 E_r^{move}&=\frac{2P_c d_r^{xy}}{3600v_c}
   +\frac{m_Rg_0(h_r^++h_r^-)}{3.6\times10^6\eta_R},\label{eq:q3-relaymove}\\
 E_r^R&=E_r^{move}+\frac{(P_h+P_{com})(30+d_r)}{3600}
       \le0.8E_R^{use}.\label{eq:q3-relayenergy}
\end{align}
其中$P_c=1.15$\,kW、$v_c=15$\,m/s、$m_R=23.5$\,kg、$\eta_R=0.72$，
$P_h+P_{com}=1.10$\,kW、$E_R^{use}=3.2$\,kWh，$g_0=9.80665\,\mathrm{m/s^2}$。
建链30\,s计入悬停与通信能耗。该式采用题设功率及前两问相同的爬升附加能耗解释。

中继机从准备开始至返航后300\,s持续被占用，任意时刻最多占用2架：
\begin{equation}
 \sum_r u_r\mathbf1_{\{t\in[c_r,R_r+300)\}}\le2\qquad(\forall t).
 \label{eq:q3-relayresource}
\end{equation}
能源组件返航后按题给SOC分段规则恢复至满电，与机体周转分别核算。
本轮候选模型对每个位置最多安排一次中继任务，每次分配不同的初始满电组件，
并施加$\sum_r u_r\le6$。这一实现保守地限制了组件复用范围；
最终方案只使用3组组件，独立核验仍记录其充电结束时刻。

\subsubsection{联合加权目标与比较条件}
沿用问题二的加权迟到定义，并将完工、能耗和架次扩展到运输与中继两类任务：
\begin{equation}
 \begin{aligned}
 W&=\sum_bw_b\max(0,C_b-d_b),\\
 C_{\max}&=\max\!\left\{\max_{p:x_p=1}(s_p+D_p),\ \max_{r:u_r=1}R_r\right\},\\
 E&=\sum_pE_px_p+\sum_rE_r^Ru_r,\qquad
 N=\sum_px_p+\sum_ru_r=N_T+N_R.
 \end{aligned}
 \label{eq:q3-objectives}
\end{equation}
联合完工时刻不包含末次返航后的充电和周转尾段，但后续任务仍须等待对应资源真正释放。
保持问题二的权重比$3:1.5:1:0.5$和固定参考尺度，目标为
\begin{equation}
 \begin{aligned}
 \min F={}&\frac12\frac{W}{429204.104}
 +\frac14\frac{C_{\max}-7528.576}{3925.283}\\
 &+\frac16\frac{E-62.634022492}{23.230078888}
 +\frac1{12}\frac{N-19}{12}.
 \end{aligned}
 \label{eq:q3-weighted}
\end{equation}
式中时间使用秒、能耗使用kWh，归一化端点与尺度在搜索过程中不调整。
参考端点不是理论下界，某项归一化值或总评分为负仍有正常的比较意义。
题给货箱优先系数$w_b$与四目标偏好权重具有不同作用。

为在已获得的按期方案之间比较通信调度，本轮主方案与对照试验均额外限定
$C_b\le d_b$对全部80箱成立，即在零期望迟到子集中搜索。
因此，当前保存方案的$W$均为0；这一性质来自附加比较条件，
不能解释为50\%的及时性权重必然保证零迟到。
硬截止、SOC、资源冲突和连续通信始终作为可行性约束，不以其他指标的改善抵消违反。

\subsection{连续通信认证与CP-SAT求解}
\subsubsection{分段航迹的连续覆盖认证}
对任一固定通信端点$\boldsymbol q$与运输直线区段
$\boldsymbol x(\lambda)=(1-\lambda)\boldsymbol a+\lambda\boldsymbol b$，$0\le\lambda\le1$，
距离函数为凸函数，最大距离不超过两个端点距离的最大值。
若该最大距离在加入10\,dB遮挡损耗后仍满足链路预算，则整个区段均可通信。
否则，进一步检查由$\boldsymbol q,\boldsymbol a,\boldsymbol b$形成的视线扫掠三角面，
逐一与经过的DEM像元柱体比较，确认整个视线面未受地形遮挡，再检验无遮挡距离上界。

当一个区段不能由同一通信端点完整认证时，对其二分并继续检查。
端点和中点检测用于发现不可覆盖的反例；连续可行性最终依赖距离界和地形几何证书。
达到20层细分或区段长度降至0.25\,m仍未取得证书时，拒绝该通信选项。
如此得到覆盖整个在途时段的连续区间，并保留每个区间的实际服务来源和链路余量下界。
多站选项同样逐区间绑定一个实际中继，其回传链路单独通过核验。

\subsubsection{候选悬停点与运输邻域增补}
候选位置先由500\,m粗网格提供，再围绕已有通信需求组的悬停位置，
在水平方向$\pm750$\,m范围内以250\,m步长补点，
结合50、100、150、200、250、275、300\,m七个离地高度进行筛选。
候选点须满足高度、往返能量、回传链路和目标运输组连续覆盖要求。
在三维坐标去重后保留65个不同位置，既比较能耗代理指标，也保留较早到位的候选。

在运输侧，采用换机型、调整访问顺序、合批、移箱和交换等操作扩充候选。
移箱及交换保留涉及各批次的完整后继组合，防止因候选截取遗漏剩余货箱。
每个新候选重新计算物理航迹和通信方案，通过筛查后再进入联合排程。
本轮每个运输候选最多访问3个服务区，并允许最多2个后备中继位置；
这两项均为实际搜索范围，不是原题对所有可行路线的限制。

\subsubsection{联合求解与独立核验}
CP-SAT同时处理运输精确覆盖、逐箱交接顺序、机体与电池占用、
中继启用、服务窗口和累计机体资源约束。
主搜索累计生成14062条物理可行运输候选，活跃池通常约350条；
邻域搜索每轮配置65\,s和8个工作线程，规范化后的固定运输结构精修配置65\,s和4个线程。
几何去重只合并同一三维位置的别名，不更改实际运输路线或中继位置。

\begin{enumerate}
 \item 从原始附件重建运输候选，冻结权重、归一化尺度及物理参数。
 \item 生成并筛选中继位置，认证每条运输航迹的直连、单站和指定站对通信选项。
 \item 建立联合整数模型，选择运输架次与通信选项，安排运输和中继执行时刻。
 \item 由完整可行方案生成配套邻域候选，交替扩充运输组合并更新排程。
 \item 固定已选运输结构精修通信关联和时序；对严格等价货箱进行不改变物理任务的标签排序。
 \item 独立复算全部交付、能耗、资源衔接与连续通信，按原固定评分比较已验证方案。
\end{enumerate}

独立核验重新读取原始附件，重算运输及中继物理量和资源时间线，
并对每个声明的通信区段重新生成连续证书。
求解返回的可行状态不能替代该核验；有限候选池及限时搜索也不构成原问题全局最优证明。

\subsection{运输与中继执行方案及可行性检验}
\subsubsection{运输路线与中继位置}
主方案执行21个运输架次，其中A、B、C型分别为8、7、6架次，完成全部80箱交付。
运输能耗为64.095969\,kWh，中继能耗为2.349304\,kWh，总能耗为66.445273\,kWh；
3个中继架次计入后，总架次数为24。
最后运输返航时刻为6559.627\,s，最后中继返航时刻为6612.284\,s，
故联合完工为1.836746\,h（110.205\,min）。
\begin{figure}[H]
 \centering
 \includegraphics[width=0.96\linewidth]{\qThreeFigDir/q3_02_routes_relays.pdf}
 \caption{推荐方案运输路线与三个实际中继悬停位置；虚线表示回传链路，不表示任意方向上的等效覆盖边界}
 \label{fig:q3-routes}
\end{figure}

运输路线见图~\ref{fig:q3-routes}，各架次配置与执行时刻见表~\ref{tab:q3-sorties}。
方案包含6个多点架次：T008访问S013$\to$S002$\to$S009，
T011访问S006$\to$S007，T015访问S003$\to$S015，T016访问S007$\to$S003，
T019访问S005$\to$S011，T021访问S011$\to$S006；其余15个架次为单点往返。
全部运输任务共包含49条实际航段，累计作业时间为40443.292\,s，
该累计量反映总作业量，与并行执行后的联合完工时间不同。
% Generated from the frozen verified Q3 solution.
\begin{table}[!htbp]
\centering
\footnotesize
\setlength{\tabcolsep}{3pt}
\renewcommand{\arraystretch}{1.10}
\caption{主方案各运输架次配置与执行时刻（均从O01出发并返回）}\label{tab:q3-sorties}
\begin{tabular}{@{}lll lrrrrr@{}}
\toprule
架次 & \shortstack{机型/\\实体机} & 电池 & 服务区访问顺序 & 箱数 & \shortstack{质量\\/kg} & \shortstack{体积\\/$\mathrm{m^3}$} & \shortstack{开始\\/min} & \shortstack{返航\\/min} \\
\midrule
T001 & A/U01 & A-01 & S006 & 2 & 22 & 0.055 & 0.00 & 21.89 \\
T002 & B/U05 & B-01 & S014 & 3 & 25 & 0.067 & 0.00 & 31.16 \\
T003 & B/U06 & B-02 & S010 & 3 & 25 & 0.067 & 0.00 & 25.92 \\
T004 & C/U07 & C-01 & S001 & 9 & 76 & 0.231 & 0.00 & 28.20 \\
T005 & C/U08 & C-02 & S001 & 6 & 78 & 0.163 & 0.00 & 24.90 \\
T006 & C/U08 & C-03 & S002 & 7 & 67 & 0.184 & 24.90 & 61.13 \\
T007 & B/U06 & B-03 & S006 & 2 & 28 & 0.054 & 25.92 & 46.04 \\
T008 & A/U01 & A-02 & S013$\to$S002$\to$S009 & 3 & 20 & 0.051 & 27.35 & 71.95 \\
T009 & A/U02 & A-03 & S013 & 2 & 22 & 0.055 & 27.85 & 54.72 \\
T010 & C/U07 & C-04 & S004 & 6 & 59 & 0.156 & 28.20 & 66.55 \\
T011 & A/U03 & A-04 & S006$\to$S007 & 3 & 20 & 0.051 & 30.75 & 66.23 \\
T012 & B/U05 & B-04 & S012 & 3 & 25 & 0.067 & 31.16 & 63.21 \\
T013 & A/U04 & A-01 & S007 & 2 & 22 & 0.055 & 41.48 & 70.27 \\
T014 & B/U06 & B-02 & S008 & 2 & 20 & 0.062 & 50.20 & 81.90 \\
T015 & A/U02 & A-05 & S003$\to$S015 & 3 & 20 & 0.051 & 54.72 & 98.11 \\
T016 & C/U08 & C-01 & S007$\to$S003 & 7 & 70 & 0.207 & 61.13 & 106.03 \\
T017 & B/U05 & B-01 & S008 & 3 & 25 & 0.067 & 63.21 & 95.91 \\
T018 & A/U03 & A-06 & S015 & 2 & 22 & 0.055 & 66.23 & 99.04 \\
T019 & C/U07 & C-02 & S005$\to$S011 & 8 & 70 & 0.196 & 66.55 & 104.97 \\
T020 & A/U01 & A-03 & S009 & 2 & 22 & 0.055 & 72.81 & 100.95 \\
T021 & B/U06 & B-03 & S011$\to$S006 & 2 & 20 & 0.062 & 81.90 & 109.33 \\
\bottomrule
\end{tabular}
\end{table}


三个悬停位置及实际服务窗口见表~\ref{tab:q3-relays}。
R001位于O01以东3.5\,km、以南0.5\,km处，离地100\,m；
R002位于以东3.5\,km、以北5.5\,km处，离地275\,m；
R003位于以西3.0\,km、以北3.0\,km处，离地275\,m。
R01实体机先后执行R001和R003，R02实体机执行R002；
3个中继架次不等于同时需要3架中继机。
% Generated from the frozen verified Q3 solution.
\begin{table}[!htbp]
\centering
\footnotesize
\setlength{\tabcolsep}{3pt}
\renewcommand{\arraystretch}{1.10}
\caption{中继任务的位置、服务窗口和返航状态}\label{tab:q3-relays}
\begin{tabular}{@{}lllrrrrrr@{}}
\toprule
任务 & \shortstack{实体机/\\组件} & \shortstack{x / y / 离地高度\\/m} & \shortstack{准备开始\\/min} & \shortstack{服务开始\\/min} & \shortstack{服务结束\\/min} & \shortstack{返航\\/min} & \shortstack{能耗\\/kWh} & \shortstack{返航SOC\\/\%} \\
\midrule
R001 & R01/RE01 & 3500 / -500 / 100 & 3.244 & 12.090 & 25.057 & 30.775 & 0.426812 & 86.662 \\
R002 & R02/RE02 & 3500 / 5500 / 275 & 26.190 & 38.729 & 60.013 & 69.650 & 0.715352 & 77.645 \\
R003 & R01/RE03 & -3000 / 3000 / 275 & 39.088 & 49.612 & 102.411 & 110.205 & 1.207141 & 62.277 \\
\bottomrule
\end{tabular}
\end{table}


\subsubsection{两类任务时序与连续通信安排}
图~\ref{fig:q3-schedule}并列展示运输和中继的执行阶段。
中继机的服务开始晚于准备开始和去程结束，只有完成建链后才能承担通信保障。
R003的有效服务延续至6144.651\,s，随后返航，于6612.284\,s结束，
比最后运输返航晚52.657\,s，决定本方案的联合完工时刻。
\begin{figure}[H]
 \centering
 \includegraphics[width=\linewidth]{\qThreeFigDir/q3_03_joint_schedule.pdf}
 \caption{21个运输架次与3个中继架次的联合执行时序}
 \label{fig:q3-schedule}
\end{figure}

图~\ref{fig:q3-communication}按实际使用来源展示各运输任务的在途通信区间。
直接链路可用的区段由G01保障，其余区段分别绑定已启用的中继任务。
T008在不同区段使用R002和R003，两站各自的服务窗口均覆盖对应需求；
该切换不涉及中继之间的多跳转发。
最终240个连续通信区间全部取得证书，未证明区间为0，
最小双向链路余量证明下界为0.079736\,dB。
这一下界基于题给静态DEM与简化传播模型，反映模型内的通信可行性。
\begin{figure}[H]
 \centering
 \includegraphics[width=\linewidth]{\qThreeFigDir/q3_04_communication.pdf}
 \caption{运输在途通信来源及中继实际服务窗口；准备和装载期间不计入在途通信条带}
 \label{fig:q3-communication}
\end{figure}

\subsubsection{逐箱交付与时限满足}
图~\ref{fig:q3-deliveries}沿用问题二的两栏逐箱表示，列出80箱实际交付时刻、期望时刻与有效硬截止。
全部货箱均按期完成，$W=0$；31个硬截止也全部满足，最小硬截止裕量为212.061\,s。
最后一个货箱于6234.755\,s完成交接，晚于该时刻的任务仅涉及返航和中继保障收尾。
逐箱所属架次、实际送达与截止时刻列于附录表~\ref{tab:q3-deliveries}。
\begin{figure}[H]
 \centering
 \includegraphics[width=\linewidth]{\qThreeFigDir/q3_05_deliveries.pdf}
 \caption{80箱交付时刻与期望时刻及有效硬截止对照}
 \label{fig:q3-deliveries}
\end{figure}

\subsubsection{机电资源占用与返航安全余量}
图~\ref{fig:q3-resources}分别展示10架实际使用机体的任务安排，以及14组运输电池和3组中继组件的周转。
同一实体机上的任务区间不重叠，运输电池复用发生在充满之后。
中继机R01在R001结束并完成300\,s周转后，才开始R003；
两趟任务分别使用RE01和RE03，不将尚未充满的组件直接用于下一架次。
图中最后返航之后的充电和周转区间表示资源恢复，不计入$C_{\max}$。
\begin{figure}[H]
 \centering
 \includegraphics[width=\linewidth]{\qThreeFigDir/q3_06_resources.pdf}
 \caption{运输与中继机体占用以及电池和组件恢复安排；虚线为联合完工时刻}
 \label{fig:q3-resources}
\end{figure}

运输无人机最低返航SOC为20.091711\%，中继机最低返航SOC为62.276848\%，均满足20\%的要求。
图~\ref{fig:q3-margins}给出硬截止裕量及逐架运输返航SOC。
最低运输SOC仅略高于规定阈值，说明该安排在题定参数下可行，
但不能由此推断其对未建模风场、能耗误差或通信环境变化具有充分裕量。
\begin{figure}[H]
 \centering
 \includegraphics[width=\linewidth]{\qThreeFigDir/q3_07_margins.pdf}
 \caption{31个硬截止货箱的时间裕量与21个运输架次的返航SOC}
 \label{fig:q3-margins}
\end{figure}

\subsubsection{完整方案核验}
独立核验共执行11884项检查，错误数为0。
核验不仅检查汇总指标，还逐一检查80箱唯一交付、站内交箱顺序、
各航段剩余载荷、机体与电池衔接、中继服务可用性及240个连续通信区间。
主要结果汇总于表~\ref{tab:q3-checks}。
% Generated from the frozen verified Q3 solution.
\begin{table}[!htbp]
\centering
\footnotesize
\setlength{\tabcolsep}{3pt}
\renewcommand{\arraystretch}{1.10}
\caption{主方案独立可行性核验结果}\label{tab:q3-checks}
\begin{tabular}{@{}>{\raggedright\arraybackslash}p{3.10cm}>{\raggedright\arraybackslash}p{5.80cm}>{\raggedright\arraybackslash}p{6.60cm}@{}}
\toprule
检查项目 & 检查条件 & 结果 \\
\midrule
货箱唯一交付 & 每箱恰好交付一次 & 80 / 80 \\
期望交付时刻 & 本轮附加零迟到条件 & 80箱全部按期，W=0 \\
有效硬截止 & 31箱，合计46条原始规则 & 全部满足；最小裕量212.061 s \\
运输返航SOC & 不低于20\% & 最低20.091711\% \\
中继返航SOC & 不低于20\% & 最低62.276848\% \\
机体与电池占用 & 同一资源不可重叠使用 & 8架运输机、14组电池均通过 \\
中继机与组件占用 & 含300 s机体周转及组件充电 & 2架中继机、3组组件均通过 \\
全在途连续通信 & 服务来源可用且取得区间证书 & 240区间通过，未证明区间0 \\
通信证明余量 & 双向接收门限包含安全裕量 & 最小下界0.079736 dB \\
独立核验汇总 & 物理、时序、通信及数据一致性 & 11884项检查，0错误 \\
\bottomrule
\end{tabular}
\end{table}


\subsection{对照分析与推荐依据}
\subsubsection{与问题二推荐方案的指标对照}
表~\ref{tab:q3-comparison}使用问题二正文采用的22架次、65.722736\,kWh、5792.178\,s方案作为基准。
本问运输架次减少1次，但加入3个中继架次后，总架次由22增加至24；
总能耗增加0.722537\,kWh，联合完工推迟820.106\,s，约13.668\,min。
运输侧能耗由65.722736\,kWh变为64.095969\,kWh，中继新增2.349304\,kWh，
两部分共同构成最终总能耗。

运输架次的减少来自重新选择运输组合后的方案差异，
不能将两份限时搜索结果理解为附加约束前后的全局最优值。
问题二的22架次是既定权重下的推荐折中方案，并非运输架次最小的证明；
问题三增加通信条件后，仍可在不同运输组合、完工时刻和能耗之间权衡。
% Generated from the frozen verified Q3 solution.
\begin{table}[!htbp]
\centering
\footnotesize
\setlength{\tabcolsep}{3pt}
\renewcommand{\arraystretch}{1.10}
\caption{各方案在固定权重与参考尺度下的指标对照}\label{tab:q3-comparison}
\begin{tabular}{@{}lrrrrrrr@{}}
\toprule
方案 & \shortstack{W\\/s} & \shortstack{运输\\架次} & \shortstack{中继\\架次} & \shortstack{总\\架次} & \shortstack{联合完工\\/s} & \shortstack{总能耗\\/kWh} & 评分F \\
\midrule
问题二推荐 & 0 & 22 & 0 & 22 & 5792.178 & 65.722736 & -0.067596986 \\
第三问主方案 & 0 & 21 & 3 & 24 & 6612.284 & 66.445273 & 0.003708106 \\
运输$\ge$22对照 & 0 & 22 & 3 & 25 & 6612.283 & 66.035704 & 0.007713987 \\
单站通信对照 & 0 & 21 & 3 & 24 & 6612.284 & 66.571180 & 0.004611438 \\
\bottomrule
\end{tabular}
\end{table}


\subsubsection{运输架次下限对照}
为检验增加运输架次后的取舍，另行加入$N_T\ge22$，中继架次不计入该下限。
其余物理规则、权重、归一化尺度、65个候选位置及全部货箱按期条件保持一致。
从原运输组的不可拆箱拆分候选出发，先保留110组配套拆分组合，
再进行联合候选搜索与固定结构精修；每个主要求解阶段配置90\,s和8个线程。
该下限仅用于本节对照，不加入正文主方案。

所得对照方案为22个运输架次与3个中继架次，总计25架次，
总能耗66.035704\,kWh，联合完工6612.283\,s，80箱仍全部按期。
独立验证11689项检查通过，并另行核对运输下限、总架次计分及固定尺度评分。
相对主方案节省0.409569\,kWh，约0.6164\%；1\,ms的完工差异不作实际提速解释。
\begin{figure}[H]
 \centering
 \includegraphics[width=\linewidth]{\qThreeFigDir/q3_08_comparison.pdf}
 \caption{问题二推荐、第三问主方案与运输下限对照；评分分解展示增加一个运输架次与节能收益的取舍}
 \label{fig:q3-comparison}
\end{figure}

固定权重与尺度下，每增加1个总架次，评分增加
$1/144\approx0.006944444$。该对照的节能使评分减少约0.002938499，
但综合评分仍由0.003708106升至0.007713987，净增量为0.004005881。
迟到及完工时刻不变时，每增加1架次，至少需要节能0.967920\,kWh才能抵消架次代价。

因此，在本次两个已验证方案中推荐21运输加3中继，22运输加3中继作为运输下限对照。
有限搜索仍可能得到更好的满足下限方案。

\subsubsection{通信策略对照与适用边界}
固定相同21条运输结构与65个候选位置，将每趟运输的后备中继位置限制为1个，
重新优化时序，得到通过验证的单站对照：总能耗66.571180\,kWh，
其中中继能耗2.475211\,kWh，联合完工6612.284\,s。
允许T008分段切换的主方案，比该单站解少消耗中继能量0.125907\,kWh。
两次限时搜索的初始方案与热启动条件不同，当前差值不能单独证明切换机制的最优增益。

当前方案使用有限的运输、站点与站对候选，并限定每站一次任务、初始满电组件分配及零迟到比较条件。
连续证书只证明静态DEM和题给传播模型内的可行性，不替代实际传播或能耗参数标定；
站点空间搜索与组件复用仍可继续扩展。

\subsection{问题三小结}
本文将地形相关链路预算、连续通信认证与中继服务时窗嵌入运输联合调度，
获得21运输加3中继的完整可行安排。全部80箱按期交付，联合完工6612.284\,s，
总能耗66.445273\,kWh，240个连续通信区间及机电资源检查均通过。
22运输加3中继对照能耗更低，但固定权重下评分较高，故保留21加3作为主方案。
上述推荐基于已验证的有限搜索结果，未证明原问题全局最优。
\par\endgroup
